%% ============================================================
%% VALIDATION RESULTS -- v2, compressed to Applied Energy style.
%% ============================================================

Table~\ref{tab:validation} reports the validation against the ex-ante gates, each metric
scored against the formulation to which it applies. The baseline prescribes supply
temperature from the heating curve and does not propagate it along pipes, so downstream
temperature metrics are structurally outside its scope and are marked as such rather than as
failures.

% CHANGED: N recomputed after reclassifying the return temperature. Two gates genuinely score
% MILP outputs: annual energy (met) and per-step closure (missed). The COP band is a
% consistency check and the two supply-temperature gates score a formulation the baseline does
% not instantiate.
Two gates score outputs of the formulation used for the cost results: the annual
delivered-energy error, which is met at $1.23$\,\%, and the per-step energy-balance closure,
which is missed at a mean of $6.4$\,\% against a $2$\,\% gate. The miss is diagnosed rather
than absorbed: it is concentrated in hours of near-zero flow, where the relative measure has a
vanishing denominator, and it does not propagate to the annual closure. The remaining entries
are a plausibility band on the heat-pump coefficient of performance, a calibration check on the
constant return temperature, and two nodal supply-temperature gates that apply to the
propagating formulation only.

\paragraph{Validated quantities.}
Two quantities test the formulation that produces every cost result, and both pass. Annual
delivered energy is reproduced to \textbf{1.23}\,\% against a $2$\,\% gate; since the network
loss is the difference between generated and delivered heat, this fixes the annual loss, which
the decomposition prices. The heat-pump coefficient of performance sits inside its plausibility
band at a mean of \textbf{2.99}; the band is a consistency check, not a measured comparison.

The source return temperature is a fixed parameter in this formulation
($63.6$\,°C, set from the design specification) rather than a solved variable, so its
\textbf{2.08}\,K mean absolute deviation from the measurement at the one sensor upstream of any
mixing valve (bias $-0.31$\,K) characterises the calibration of that parameter and the scatter
of the measured series about it. It is reported as such and not counted among the validated
outputs. Return-temperature propagation is available in the nonlinear formulation, where the
temperature is a free variable; the present results do not depend on it.

% CHANGED: units. Trunk drop, far-end error, bias and the valve offset are temperature
% DIFFERENCES and belong in K, not degC; absolute levels (63.6 degC) correctly stay in degC.
% CHANGED: the trunk-drop gate is named as mis-specified rather than merely missed -- a gate on
% a quantity the formulation does not produce cannot be satisfied by any calibration.
\paragraph{Temperature field: limitation.}
The supply-temperature field belongs to the propagating formulation and is not quantitatively
validatable on this network. The baseline does not propagate supply temperature, so its trunk
drop is zero by construction against a measured $8.2$\,K. This gate was mis-specified: no
calibration of the baseline can satisfy it, because the quantity it scores is not one the
formulation produces. It is therefore reported as inapplicable to the baseline rather than as a
failure, as are the far-end and corridor gates. On the propagating formulation the far-end error
is $9.21$\,K and is almost entirely bias (mean absolute error and bias coincide to four
decimals), matching the $5$--$15$\,K offset of a three-way mixing valve rather than model error;
the corridor comparison shows the same signature (bias $-7.73$\,K). A model resolved more finely
than its metering cannot have the extra resolution verified, so the conclusions are anchored to
the annual network loss, which both formulations agree on and the measurements resolve. The
far-end node over the full annual record is reported rather than a mean over selected nodes, as
the more conservative choice; the all-node mean of $1.68$\,K lies between the calibration- and
validation-node group means (Appendix~\ref{app:per_node_mae}).

% CHANGED: added the maximum of the per-step residual, so the table's 99 % is not read as model
% failure. CHANGED: deleted the closing block, which repeated the head paragraph verbatim and
% contradicted its gate count (four vs. two gates scoring the MILP formulation).
\paragraph{Flow and energy balance.}
The source flow shows a mean absolute percentage error of \textbf{33.8}\,\%, governed by the
denominator rather than the model: the absolute error is near-constant at $12$--$13$\,m$^3$/h,
so the same error is $46$\,\% of mean flow below a quarter of peak load but $13$\,\% at half to
three-quarters load. Low-load hours are $60$\,\% of the year but carry little energy, so the
large relative error is compatible with the $1.23$\,\% annual energy match. First differences
separate a level error from a variation error, since a fixed offset cancels: the model tracks
flow level at $r=\textbf{0.91}$ and day-to-day change at $\textbf{0.80}$, and demand at $0.93$
and $0.89$. The per-step energy balance averages $6.4$\,\%, dominated by the same low-flow
hours, and reaches $99$\,\% in the extreme, in hours where flow is essentially zero: the
relative measure degenerates as its denominator vanishes, while the absolute residual in those
hours remains negligible. The annual closure of $1.23$\,\% is correspondingly unaffected and
well inside its gate.

\paragraph{Post-upgrade assets.}
The heat pump, electrode boiler and storage were installed after the measurement record, so
their dispatch is verified against necessary conditions rather than measurements: coefficient
of performance in band, storage state of charge within limits, and no simultaneous charge and
discharge. Two structural checks close the section: the nonlinear reference with its quadratic
constraints deactivated reproduces the baseline objective to better than $0.01$\,\%, and the
relaxation property $\mathrm{cost}(\CP) \le \mathrm{cost}(\Lone)$ holds on every one of the 135
synthetic instances.